\section{Introduction}
\label{sec:introduction}

The digitalization and electrification of transportation are driving a rapid convergence between transmission networks, distribution feeders, and large fleets of software managed electric vehicles (EVs) and distributed energy resources (DERs). Modern grid operations increasingly depend on cloud based management platforms, supervisory control and data acquisition (SCADA), and programmatic interfaces that schedule EV charging and DER setpoints across thousands of devices. While this integration improves observability and flexibility, it also exposes new logic level attack surfaces: adversaries who compromise EV or DER management platforms can coordinate load altering attacks that propagate from distribution feeders into the bulk transmission system, potentially causing line overloads, voltage instability, and large scale load shedding \cite{morris2013icssecurity,cintuglu2017testbeds,yohanandhan2022review,laaSurvey2024}.

To assess grid resilience, the community has developed a rich ecosystem of testbeds, co simulation frameworks, and intrusion detection systems (IDS) for industrial control systems and smart grids \cite{cintuglu2017testbeds,hardy2024helics,yohanandhan2022review}. However, most evaluation methodologies still rely on static datasets or a small number of pre scripted attack scenarios. Widely used traces such as SWaT, WADI, EPIC, and HAI capture fixed attack trajectories under specific operating conditions \cite{swatdataset,wadidataset,epicdataset,haidataset}; many smart grid studies likewise use hand crafted attack scripts that do not change during an experiment \cite{le2020gridattacksim}. We refer to the resulting discrepancy between a defense's reported performance against such static benchmarks and its robustness against reactive adversaries as the \emph{evaluation validity gap}. Static traces and scripts cannot capture the action reaction dynamics of adaptive attackers and can therefore produce misleading confidence in deployed defenses \cite{sahani2023smartgridids}.

Concurrently, advances in large language models (LLMs) and tool using agents have transformed offensive capabilities in IT and industrial control domains. Systems such as PentestGPT, AutoAttacker, and AttackLLM demonstrate that LLM based agents can decompose complex tasks, read documentation, orchestrate external tools, and refine strategies based on feedback \cite{deng2024pentestgpt,autoattacker,attackllm,fang2024llmagents}. In this work, we view LLMs as scalable \emph{cognitive proxies} for human red teams: they provide semantic reasoning and decision making over heterogeneous artifacts (protocols, APIs, logs) and can be embedded into automated evaluation pipelines.

Applying LLM agents to cyber physical systems introduces a key challenge: traditional LLM based red teaming frameworks are largely IT focused and are not aware of physical constraints. In power systems, hallucinated commands---such as impossible voltages, infeasible power injections, or violating device ratings---are not merely incorrect but physically meaningless. To obtain meaningful evaluations, any LLM based attacker for smart grids must therefore be grounded in a \emph{physically constrained primitive interface} that exposes only valid, well typed actions over the grid state, such as EV capacity adjustments within configured bounds.

This paper introduces LLM-GridEval, an evaluation framework that embeds an adaptive, tool using LLM attacker inside a HELICS based transmission and distribution co simulation \cite{hardy2024helics}. LLM-GridEval connects the attacker to the smart grid through a schema constrained catalog of attack primitives (e.g., \texttt{get\_grid\_status()}, \texttt{set\_ev\_capacity}) that are enforced by a sandboxed interface layer. The framework is explicitly designed to measure and analyze the evaluation validity gap by comparing defenses under static, scripted attackers and under adaptive, observation driven policies.

The contributions of this paper are threefold:
\begin{itemize}
	\item \textbf{Theoretical:} We formalize defense evaluation against adaptive attackers in cyber physical power systems by modeling the attacker grid interaction as a partially observable process and by defining the evaluation validity gap between static and adaptive attack processes.
	\item \textbf{Methodological:} We design LLM-GridEval, a framework that combines HELICS based T\&D co simulation with a schema constrained attack primitive catalog and a sandboxed interface layer suitable for LLM based and algorithmic agents.
	\item \textbf{Empirical:} We instantiate an EV capacity load altering case study in which attacker issued EV setpoint changes, applied through the constrained primitive interface in a realistic T\&D co simulation, drive a distribution feeder into overload and solver instability, illustrating how adaptive attackers can expose vulnerabilities that static baselines fail to reveal.
\end{itemize}
